\documentclass[11pt, a4paper]{article}

% --- UNIVERSAL PREAMBLE BLOCK ---
\usepackage[a4paper, top=2.5cm, bottom=2.5cm, left=2cm, right=2cm]{geometry}
\usepackage{fontspec}

\usepackage[italian, bidi=basic, provide=*]{babel}

\babelprovide[import, onchar=ids fonts]{italian}
\babelprovide[import, onchar=ids fonts]{english}

% Set default/Latin font to Sans Serif in the main (rm) slot
\babelfont{rm}{Noto Sans}

% Add because main language is not English
\usepackage{enumitem}
\setlist[itemize]{label=-}

% --- ALTRI PACCHETTI ---
\usepackage{hyperref}
\usepackage{listings}
\usepackage{xcolor}
\usepackage{titlesec}
\usepackage{fancyhdr}

% Configurazione Colori e Stile per il Codice
\definecolor{codegray}{rgb}{0.5,0.5,0.5}
\definecolor{codepurple}{rgb}{0.58,0,0.82}
\definecolor{backcolour}{rgb}{0.95,0.95,0.92}
\definecolor{uniblue}{rgb}{0.1, 0.2, 0.6}

\lstdefinestyle{mystyle}{
    backgroundcolor=\color{backcolour},   
    commentstyle=\color{codegray},
    keywordstyle=\color{uniblue}\bfseries,
    numberstyle=\tiny\color{codegray},
    stringstyle=\color{codepurple},
    basicstyle=\ttfamily\small,
    breakatwhitespace=false,         
    breaklines=true,                 
    captionpos=b,                    
    keepspaces=true,                 
    numbers=left,                    
    numbersep=5pt,                  
    showspaces=false,                
    showstringspaces=false,
    showtabs=false,                  
    tabsize=4
}
\lstset{style=mystyle}

% Configurazione Link
\hypersetup{
    colorlinks=true,
    linkcolor=uniblue,
    filecolor=magenta,      
    urlcolor=uniblue,
    pdftitle={Documentazione Progetto LSO - Tris},
    pdfauthor={Luca Pagan}
}

\begin{document}

% ---------------------------------------------------------
% FRONTESPIZIO
% ---------------------------------------------------------
\begin{titlepage}
    \centering
    \vspace*{1cm}
    
    {\Large \textbf{UNIVERSITÀ DEGLI STUDI}}\\[0.5cm]
    {\large Dipartimento di Ingegneria Elettrica e delle Tecnologie dell'Informazione}\\[2cm]
    
    {\Large Corso di Laboratorio di Sistemi Operativi}\\[2.5cm]
    
    \rule{\textwidth}{0.4pt}\\[0.5cm]
    {\huge \textbf{Documentazione del Progetto\\[0.5cm] Tic-Tac-Toe (Tris) Client-Server}}\\[0.4cm]
    \rule{\textwidth}{0.4pt}\\[3cm]
    
    \begin{flushleft}
        \Large
        \textbf{Studenti:}\\
        Luca Pagano\\
        \normalsize Matricola: N86004788\\
        \Large
        Danilo Scala\\
        \normalsize Matricola: N86004474\\
    \end{flushleft}
    
    \vfill
    
    {\large Anno Accademico 2024-2025}\\[1cm]
    
\end{titlepage}

\newpage
\tableofcontents
\newpage

% ---------------------------------------------------------
% SEZIONE 1: INTRODUZIONE
% ---------------------------------------------------------
\section{Introduzione e Requisiti}

Il presente documento illustra le scelte progettuali e implementative relative alla realizzazione di un'applicazione client-server dedicata alla gestione di partite del gioco del Tris in modalità multi-client. Il progetto è stato sviluppato come elaborato finale per il corso di Laboratorio di Sistemi Operativi.

L'architettura del sistema prevede un \textbf{server centrale} scritto in linguaggio C, che si occupa di coordinare le connessioni e la logica di gioco, e uno o più \textbf{client} (sviluppati tipicamente con interfaccia grafica in Java/JavaFX o da terminale) che permettono agli utenti di interagire con il sistema.

\subsection{Regole e Stati della Partita}
Un utente connesso può creare nuove partite o partecipare a partite esistenti. Il ciclo di vita di una partita è gestito attraverso una macchina a stati:
\begin{itemize}
    \item \textbf{IN ATTESA}: La partita è stata creata da un utente ed è in attesa che un secondo giocatore si unisca.
    \item \textbf{IN CORSO}: Un secondo giocatore ha richiesto di partecipare, il creatore ha accettato la sfida e i giocatori si stanno alternando nei turni.
    \item \textbf{TERMINATA}: La partita si è conclusa con una vittoria, una sconfitta o un pareggio.
\end{itemize}

In caso di pareggio, il sistema prevede la possibilità di richiedere una \textbf{rivincita}. Se entrambi i client accettano, la tabella di gioco viene resettata e lo stato torna \textit{IN CORSO}.

\newpage
% ---------------------------------------------------------
% SEZIONE 2: DESIGN DEL SISTEMA
% ---------------------------------------------------------
\section{Design del Sistema}

Questa sezione descrive nel dettaglio le scelte progettuali relative all'architettura distribuita, alla gestione dell'ambiente di esecuzione e agli strumenti utilizzati per garantire scalabilità, manutenibilità e facilità di sviluppo dell'intero ecosistema client-server.

\subsection{Comunicazione di Rete}
La comunicazione tra client e server avviene esclusivamente tramite \textbf{socket TCP} (Transmission Control Protocol). Questa scelta architetturale è stata ponderata per garantire due aspetti fondamentali:
\begin{itemize}
    \item \textbf{Affidabilità}: il protocollo TCP assicura una consegna ordinata e senza perdite dei messaggi, condizione fondamentale e imprescindibile per mantenere la coerenza dello stato del gioco e del tabellone tra i vari partecipanti.
    \item \textbf{Connessione persistente}: una volta stabilita la connessione iniziale, il canale socket resta attivo per l'intera durata della sessione di gioco, evitando l'overhead di riconnessione ad ogni singolo messaggio scambiato.
\end{itemize}

\subsection{Containerizzazione e Deployment}
Per semplificare il processo di deployment e garantire l'assoluta portabilità del software, il sistema è stato interamente containerizzato utilizzando \textbf{Docker}. Sono stati definiti e creati due container distinti e isolati:
\begin{itemize}
    \item \textbf{Container Server}: predispone un ambiente minimale contenente il server scritto in C e tutte le librerie di sistema necessarie.
    \item \textbf{Container Client}: fornisce un ambiente preconfigurato in Java per l'esecuzione fluida dell'applicativo client.
\end{itemize}
In questo modo è possibile eseguire l'intero sistema su qualunque macchina host azzerando i conflitti di dipendenze e assicurando una maggiore riproducibilità degli esperimenti.

\subsection{Gestione del Codice, Configurazione e Build}
Il ciclo di vita del software è gestito attraverso rigorose pratiche di automazione:
\begin{itemize}
    \item \textbf{Automazione della Build}: La compilazione e l'esecuzione dei due moduli principali sono state automatizzate per ottimizzare i tempi di sviluppo:
    \begin{itemize}
        \item \textbf{Makefile}: utilizzato per automatizzare la compilazione del backend in C.
        \item \textbf{Maven}: impiegato per la gestione delle dipendenze e la build del frontend in Java.
    \end{itemize}
\end{itemize}

\newpage
% ---------------------------------------------------------
% SEZIONE 3: BACKEND (SERVER)
% ---------------------------------------------------------
\section{Architettura del Backend (Server)}

Il backend rappresenta il cuore del sistema e gestisce l'integrità del gioco e la concorrenza tra i client. All'avvio, il server crea un socket passivo (\texttt{listen}) su una porta predefinita.

\subsection{Gestione della Concorrenza}
Per consentire a molteplici utenti di giocare contemporaneamente senza bloccare il server, è stato adottato un modello \textbf{Thread-per-Client}. 
Per ogni connessione accettata dalla \texttt{accept()}, il server delega la gestione del client ad un thread dedicato creato tramite \texttt{pthread\_create}.

\subsection{Strutture Dati Condivise}
Il server deve mantenere lo stato globale del sistema, che comprende l'elenco degli utenti connessi e l'elenco delle partite.

\begin{lstlisting}[language=C, caption=Struttura Dati Client]
struct client {
    unsigned int clientfd;        // File descriptor del socket
    char* name;                   // Username del giocatore
    pthread_mutex_t sock_mutex;   // Mutex per accesso esclusivo al socket
};
\end{lstlisting}

Il \texttt{sock\_mutex} è fondamentale: previene il fenomeno dell'\textit{interleaving}. Quando il thread del server deve inviare una sequenza logica di messaggi al client (es. aggiornamento del tabellone e cambio turno), il mutex impedisce ad altri thread di scrivere sul medesimo socket contemporaneamente.

\begin{lstlisting}[language=C, caption=Struttura Dati Game]
struct game {
    unsigned short int id_game;
    struct client *client1;
    struct client *client2;
    unsigned int **table;         // Matrice 3x3 per il tabellone
    game_status status;           // ATTESA, IN_CORSO, TERMINATA
    unsigned int turn;            // Indica a chi tocca muovere
};
\end{lstlisting}

\subsection{La libreria DynamicArray e Thread-Safety}
Le liste dei client e dei game sono gestite tramite una struttura personalizzata \texttt{DynamicArray}. Poiché l'array dinamico è una risorsa condivisa tra tutti i thread (che possono aggiungere o rimuovere client e partite in qualsiasi momento), esso include un \texttt{pthread\_mutex\_t} interno.

Ogni operazione di inserimento (\texttt{da\_push}), lettura (\texttt{da\_get\_at}) o rimozione (\texttt{da\_remove\_at}) acquisisce il lock prima di manipolare l'array, garantendo totale \textbf{thread-safety}. Inoltre, l'array è capace di ridimensionarsi automaticamente (\textit{shrink}) quando la memoria inutilizzata supera una certa soglia.

\subsection{Generazione degli ID in tempo \texorpdfstring{$O(1)$}{O(1)}}
\begin{sloppypar}
Al fine di identificare in modo efficiente le partite, il sistema sfrutta gli indici del \texttt{DynamicArray}. La funzione \texttt{generate\_id()} restituisce il primo indice libero. Quando termina una partita, \texttt{release\_id()} segna quell'indice come riutilizzabile. Questo garantisce un accesso al match in tempo costante $O(1)$ usando il suo ID, evitando di iterare sull'intero array.
\end{sloppypar}

\subsection{Gestione dei Comandi (Protocollo)}
Il protocollo è \textit{command-based}. I messaggi iniziano con un \texttt{enum command} (intero) che indica l'azione richiesta.

\begin{lstlisting}[language=C, caption=Enum dei Comandi]
typedef enum {
    SERVER_OK = 0,
    CREATE_GAME,
    GET_QUEUE,
    START_GAME,
    DISCONNECT,
    ACCEPT,
    REJECT,
    MAKE_MOVE
} command;
\end{lstlisting}

Il thread dedicato al client legge ciclicamente il socket. Quando riceve \texttt{MAKE\_MOVE}, per esempio, aggiorna la matrice del gioco, verifica le condizioni di vittoria, e invia un broadcast (protetto da mutex) ai due socket coinvolti nella partita per aggiornare l'interfaccia. Oltre all'enumerazione di base, i payload dei messaggi vengono formattati e scambiati strutturandoli nel formato \textbf{JSON} (tramite l'integrazione della libreria esterna \texttt{cJSON}). Questa scelta garantisce un'elevata flessibilità e facilità di parsing per il trasferimento di dati strutturati, come le coordinate delle mosse o i nomi utente, rendendo la comunicazione tra C e Java robusta ed estensibile.

\newpage
% ---------------------------------------------------------
% SEZIONE 4: FRONTEND (CLIENT)
% ---------------------------------------------------------
\section{Architettura del Frontend (Client)}

Il client fornisce all'utente finale uno strumento per interagire col server. Se sviluppato in Java con \textbf{JavaFX}, esso adotta il pattern architetturale \textit{Model-View-Controller (MVC)}.

\subsection{Componenti Logici}
\begin{itemize}
    \item \textbf{Views (FXML)}: Definiscono la struttura visiva. Troviamo tipicamente una \texttt{HomeView} per la gestione delle lobby e una \texttt{GameView} che renderizza la griglia 3x3 del Tris.
    \item \textbf{Controllers}: Classi Java associate alle View. Catturano gli eventi dell'interfaccia (es. click su un pulsante del tabellone) e invocano i metodi del Client di rete.
    \item \textbf{Network Handler (GameClient)}: È il modulo responsabile dell'astrazione del livello di rete. Gestisce il socket TCP lato client.
\end{itemize}

\subsection{Gestione Asincrona della Rete}
Un aspetto critico del client è la separazione tra il thread dell'Interfaccia Grafica (UI Thread) e l'I/O di rete. Le operazioni di rete in lettura sono bloccanti. Se il client leggesse dal socket nel thread della UI, l'applicazione si bloccherebbe in attesa dei messaggi del server ("Application Not Responding").

Per ovviare a ciò, il \texttt{GameClient} avvia un thread demone parallelo (\texttt{startReader()}) dedicato esclusivamente all'ascolto dei messaggi provenienti dal server. Quando arriva un evento (es. "Avversario ha mosso"), il thread di rete utilizza primitive come \texttt{Platform.runLater()} (in JavaFX) per aggiornare la UI in modo thread-safe.

\newpage
% ---------------------------------------------------------
% SEZIONE 5: LOGICA DI GIOCO
% ---------------------------------------------------------
\section{Logica e Controllo del Tris}

La validazione delle mosse e la determinazione dell'esito della partita sono gestite dal Server, per evitare che client malintenzionati possano inviare mosse fraudolente (Server-Authoritative model).

\subsection{Ciclo del Turno}
\begin{enumerate}
    \item Il client di turno invia il comando \texttt{MAKE\_MOVE} con le coordinate $(x, y)$.
    \item Il server acquisisce il lock della partita, verifica che la cella $(x, y)$ sia vuota e che sia effettivamente il turno di quel client.
    \item La cella viene marcata con l'ID o il simbolo del giocatore.
    \item Vengono richiamate le funzioni di controllo: \texttt{check\_row()}, \texttt{check\_col()} e \texttt{check\_diag()}.
    \item Se viene rilevata una vittoria, il server invia un pacchetto \texttt{GAME\_OVER} ad entrambi, seguito dallo stato di vittoria o sconfitta.
    \item In caso contrario, verifica se il tabellone è pieno (\texttt{isBoardFull()}). Se sì, decreta il pareggio.
    \item Altrimenti, il turno viene passato all'avversario e viene inviato il comando di aggiornamento ad entrambi i client.
\end{enumerate}

\subsection{Gestione delle Disconnessioni}
Se un utente chiude brutalmente il client (crash o perdita di rete), la funzione \texttt{recv()} sul server restituisce 0 o un errore. Il thread del server rileva l'evento:
\begin{itemize}
    \item Libera il file descriptor.
    \item Rimuove il client dalla lista dei client attivi.
    \item Se il client era in una partita \textit{IN CORSO}, la partita viene terminata e viene notificata all'avversario la "Vittoria per abbandono".
\end{itemize}

\newpage
% ---------------------------------------------------------
% SEZIONE 6: CONCLUSIONI
% ---------------------------------------------------------
\section{Conclusioni}
Il progetto ha permesso di applicare concretamente i concetti fondamentali dei sistemi operativi: gestione della concorrenza, sincronizzazione tra thread, utilizzo di meccanismi di mutua esclusione (mutex) e comunicazione IPC (Inter-Process Communication) tramite socket di rete. 

L'architettura proposta garantisce stabilità, isolamento delle risorse e una gestione efficiente della memoria grazie alla corretta deallocazione e al riciclo degli identificatori logici. Durante lo sviluppo, al fine di scongiurare la presenza di memory leak o corruzioni di memoria nel backend, il codice C è stato rigorosamente analizzato e compilato in fase di test con l'ausilio di \texttt{AddressSanitizer} (\texttt{-fsanitize=address}).

Inoltre, la robustezza del codice è stata validata attraverso l'implementazione di una suite di Unit Test. La solidità della logica del gioco (ricerca delle condizioni di vittoria, validazione delle mosse) e l'integrità delle strutture dati personalizzate (come l'array dinamico thread-safe) sono state testate sia tramite asserzioni standard nel codice C, sia avvalendosi del framework standard industriale \textbf{Google Test} (\texttt{gtest}) in C++, garantendo in tal modo una copertura affidabile per le successive iterazioni di sviluppo.

\end{document}